% W5i77_HardProblems.tex
% DFD 20-Agent Research Programme --- Iteration 77
% Hard Problems, Honest Negatives, and Open Questions

\documentclass[11pt,a4paper]{article}
\usepackage[utf8]{inputenc}
\usepackage{amsmath,amssymb}
\usepackage{booktabs}
\usepackage{geometry}
\usepackage{xcolor}
\usepackage{enumitem}
\geometry{margin=2.5cm}

\title{W5i77: Hard Problems and Honest Negatives\\[6pt]
\large Zenodo Reproducibility Gaps $\cdot$ Bispectrum Interpretation $\cdot$ Portfolio Bottleneck}
\author{DFD 20-Agent Research Programme}
\date{2026-03-24}

\begin{document}
\maketitle

\section{Hard Problem 1: Zenodo Reproducibility on Windows}

The Zenodo archive has been tested on macOS and Ubuntu but not Windows. Known issues:
\begin{itemize}[nosep]
\item Julia DFD packages use POSIX path separators in two scripts (DFDBolt.jl line 142, DFDGravity.jl line 87). Fix: use \texttt{joinpath()} throughout.
\item Mathematica notebooks reference Unix-style temp directories. Fix: use \texttt{\$TemporaryDirectory}.
\item Python scripts are cross-platform (no issues expected).
\end{itemize}

\textbf{Severity}: LOW. Fixes are trivial (estimated 2 hours). Not a scientific issue.

\textbf{Mitigation}: Fix path separators in i78. Test on Windows VM before Zenodo upload.

\section{Hard Problem 2: Bispectrum Enhancement Interpretation}

The N-body bispectrum shows $B^{\rm DFD}/B^{\Lambda{\rm CDM}} = 1.018 \pm 0.006$ at $k = 0.2\,h$/Mpc (equilateral). Two interpretive questions:

\begin{enumerate}[nosep]
\item \textbf{Is the enhancement from $\psi$-coupling or from modified halo profiles?} The squeezed configuration enhancement ($1.022$) is larger than equilateral ($1.018$), suggesting a mode-coupling origin (large-scale $\psi$ modulates small-scale growth). This is consistent with the DFD $\psi$-coupling mechanism. However, a definitive separation requires a perturbative bispectrum calculation (not yet done).

\item \textbf{Is $1.8\%$ detectable?} Current galaxy survey bispectrum measurements have $\sim 5$--$10\%$ errors (BOSS). Euclid will achieve $\sim 2$--$3\%$ (Yankelevich \& Porciani 2019). So $1.8\%$ is marginally detectable with Euclid. Combined with power spectrum: potentially significant.
\end{enumerate}

\textbf{Honest negative}: The bispectrum alone is probably not a discovery channel for DFD. It is a supporting probe, not a primary one.

\section{Hard Problem 3: Portfolio Bottleneck}

With 8 submission-ready papers + 1 submitted, the programme faces a practical bottleneck:
\begin{itemize}[nosep]
\item Submitting all 8 simultaneously risks dilution (referees, community attention).
\item Sequential submission over 6--12 months is more strategic but slower.
\item Optimal strategy: submit 2--3 highest-impact papers first, then follow up.
\end{itemize}

Recommended submission order:
\begin{enumerate}[nosep]
\item v3.3 (Physics Reports) + $\alpha$ PRL letter --- simultaneous, flagship pair.
\item $E_G$ paper (PRD) + No-screening paper (JCAP) --- gravity sector pair.
\item DESI paper (JCAP) + Euclid pre-registration (JCAP) --- data confrontation pair.
\item Review paper + Software/JOSS --- supporting material.
\end{enumerate}

\textbf{Not a scientific problem}: this is a strategic/logistical issue. All papers are scientifically ready.

\section{Hard Problem 4: N-Body CPU Budget}

Current allocation: 500k CPU-hours. Consumed: 320k. Remaining: 180k. Needed for completion:
\begin{itemize}[nosep]
\item Complete $z = 0$--$2$ evolution: $\sim 80$k hours. \checkmark\ (within budget)
\item $L = 1000\,h^{-1}$Mpc BAO box: $\sim 800$k hours. \textbf{Not within current budget.}
\item Cosmology variation suite (5 $\Omega_m$ values): $\sim 400$k hours. Not within budget.
\end{itemize}

The primary N-body paper (power spectrum, halo MF, voids, bispectrum) can be completed within budget. BAO-scale and cosmology variation require a Phase 2 HPC proposal.

\section{Hard Problem 5: $E_G$ and $S_8$ Tension Landscape}

The DFD $E_G$ paper predicts $E_G > E_G^{\rm GR}$. But DFD also predicts $S_8 = 0.801$, close to Planck ($0.832$) but above weak lensing measurements ($0.76$--$0.79$). If the $S_8$ tension resolves in favour of CMB ($S_8 \sim 0.83$), this is fine. If it resolves in favour of lensing ($S_8 \sim 0.77$), DFD has a $\sim 2\sigma$ tension in $S_8$.

\textbf{Honest negative}: DFD does not naturally resolve the $S_8$ tension. It reduces it (from $3\sigma$ to $\sim 2\sigma$) but does not eliminate it. The DFD $\psi$-screen affects $\Sigma$ (lensing) more than $\mu$ (clustering), producing the $E_G$ signature but not a large $S_8$ shift.

\textbf{Not a falsification risk}: the $S_8$ tension is not definitively measured. DFD is within the current error bars.

\section{Open Questions from i77}

\begin{enumerate}[nosep]
\item What is the optimal submission timeline for the 8-paper portfolio?
\item Can the N-body bispectrum enhancement be decomposed into $\psi$-coupling vs halo profile contributions analytically?
\item Is a Phase 2 HPC proposal competitive for $\sim 1$M CPU-hours?
\item Should the Zenodo archive include raw N-body snapshots (large) or only post-processed statistics?
\item What is the expected referee turnaround for Physics Reports (v3.3)?
\end{enumerate}

\section{Summary}

5 hard problems identified. 0 are RED risks. 1 is a trivial technical fix (Windows paths). 1 is an interpretive nuance (bispectrum). 1 is logistical (portfolio timing). 1 is resource-limited (HPC budget). 1 is an honest limitation ($S_8$). Programme integrity: \textbf{MAINTAINED}.

\end{document}
